\documentclass[11pt,a4paper]{article}

% --- Encoding and fonts ---
\usepackage[T1]{fontenc}
\usepackage[utf8]{inputenc}
\usepackage{mathpazo}          % Palatino text + Pazo math
\linespread{1.05}              % Palatino benefits from slightly more leading
\usepackage{microtype}         % microtypographic improvements (hyphenation, spacing)

% --- Core math and symbols ---
\usepackage{amsmath}
\usepackage{amsfonts}
\usepackage{amssymb}

% --- Graphics and colour ---
\usepackage{graphicx}
\usepackage{xcolor}

% --- Tables ---
\usepackage{booktabs}          % \toprule, \midrule, \bottomrule
\usepackage{tabularx}
\usepackage{float}

% --- Lists ---
\usepackage{enumitem}
\setlist{itemsep=2pt, topsep=4pt, parsep=0pt}

% --- Layout ---
\usepackage{geometry}
\geometry{a4paper, top=1.2in, bottom=1.2in,
          left=1.15in, right=1.15in, headheight=14pt}

% --- Headers and footers ---
\usepackage{fancyhdr}
\pagestyle{fancy}
\fancyhf{}
\fancyhead[L]{\small\itshape\nouppercase{\leftmark}}
\fancyhead[R]{\small\itshape StarkZee Manual}
\fancyfoot[C]{\small\thepage}
\renewcommand{\headrulewidth}{0.4pt}
\renewcommand{\footrulewidth}{0pt}

% --- Section heading styles ---
\usepackage{titlesec}
\titleformat{\section}
  {\large\bfseries}{\thesection}{0.8em}{}
  [\vspace{2pt}\rule{\linewidth}{0.5pt}\vspace{2pt}]
\titleformat{\subsection}
  {\normalsize\bfseries}{\thesubsection}{0.8em}{}
\titleformat{\subsubsection}
  {\normalsize\itshape}{\thesubsubsection}{0.8em}{}
\titlespacing*{\section}    {0pt}{18pt plus 2pt minus 1pt}{8pt}
\titlespacing*{\subsection} {0pt}{12pt plus 1pt minus 1pt}{4pt}
\titlespacing*{\subsubsection}{0pt}{8pt plus 1pt}{2pt}
\setlength{\emergencystretch}{3em}  % fallback for lines with long \texttt{} identifiers

% --- Hyperlinks ---
\usepackage[colorlinks=true,
            linkcolor={[HTML]{1A3A6B}},
            citecolor={[HTML]{1A6B3A}},
            urlcolor={[HTML]{1A3A6B}},
            pdftitle={The Stark-Zeeman Line-Shape Code},
            pdfauthor={G. Ronchi}]{hyperref}

% --- Code listings ---
\usepackage{listings}

% --- Title ---
\title{%
  \Large\bfseries The Stark-Zeeman Line-Shape Code\\[6pt]
  \large\mdseries\itshape User Manual and Theoretical Guide}
\author{\textbf{G.\ Ronchi}\\[4pt]
  \small\itshape Open-Source Coupled Stark-Zeeman Plasma Line-Shape Project}
\date{June 2026}

\begin{document}

\maketitle
\thispagestyle{empty}

\tableofcontents
\newpage

\section{Introduction}
Modeling the spectral line shapes of multi-electron ions in hot, dense plasmas is a challenging problem in plasma physics, astrophysics, and laboratory nuclear fusion diagnostics. The presence of strong external magnetic fields ($\vec{B}$) coupled with the stochastic electric microfields ($\vec{F}$) produced by moving ions and electrons in the plasma leads to a complex, simultaneous splitting and broadening known as the Stark-Zeeman (SZ) effect.

The \textbf{Stark Zeeman Line-Shape} (\textbf{StarkZee}) code is an open-source Python implementation that couples atomic physics calculations in arbitrary magnetic fields with plasma line-broadening theories under quasi-static electric microfields and electron impact. The library incorporates ion dynamics and electric field fluctuations via the Frequency Fluctuation Model (FFM).

This manual details the theoretical background, governing equations, code structure, physical simplifications, and high-performance NumPy vectorization optimizations implemented in the package.

\section{Theoretical Background and Governing Equations}

\subsection{Liouville Space Representation}
The spectral line-shape intensity $I(\omega)$ of emitted radiation is related to the Fourier transform of the dipole autocorrelation function $C(t)$ by \cite{baranger}:
\begin{equation}
\label{eq:Iomega}
I(\omega) = \frac{1}{\pi} \operatorname{Re} \int_0^\infty C(t)\, e^{i\omega t}\, dt
\end{equation}
In Liouville space notation, the autocorrelation function is a trace over the quantum emitter states:
\begin{equation}
\label{eq:Ct}
C(t) = \langle\langle \vec{d}^{\,*} | U(t) | \vec{d}\,\rho_0 \rangle\rangle
\end{equation}
where $\vec{d}$ is the electric transition dipole operator, $\rho_0$ is the equilibrium density matrix of the initial manifold, and $U(t)$ is the time-evolution propagator.

In the anti-symmetric subspace (ground-to-excited manifold transitions in the no-quenching approximation), the Liouvillian operator $L$ is constructed from the excited-state Hamiltonian $H_u$ and ground-state Hamiltonian $H_l$:
\begin{equation}
\label{eq:liouvillian}
L = \frac{1}{\hbar} \bigl(H_u \otimes \hat{I}^d - \hat{I} \otimes H_l^d\bigr)
\end{equation}
The diagonal elements of $L$ give the field-free transition frequencies $\omega_{ij} = (E_i^u - E_j^l)/\hbar$; off-diagonal elements encode the Stark and Zeeman couplings between dressed states.

\paragraph{From propagator to resolvent.}
The time-evolution propagator is generated by $L$:
\begin{equation}
\label{eq:propagator}
U(t) = e^{-iLt}
\end{equation}
Substituting into Eq.~(\ref{eq:Ct}) gives the explicit correlation function
\begin{equation}
\label{eq:Ct_explicit}
C(t) = \langle\langle \vec{d}^{\,*} | e^{-iLt} | \vec{d}\,\rho_0 \rangle\rangle
\end{equation}
Inserting this into Eq.~(\ref{eq:Iomega}) and evaluating the one-sided Fourier transform converts the time-domain exponential into a frequency-domain resolvent:
\begin{equation}
\label{eq:resolvent_bare}
I(\omega) = \frac{1}{\pi}\operatorname{Re}\,
  \langle\langle \vec{d}^{\,*} \big|\,(\omega\,\hat{I} - L)^{-1}\,\big| \vec{d}\,\rho_0 \rangle\rangle
\end{equation}
Fast electron collisions are encoded as an additional imaginary damping operator
$\Phi(\omega)$ in the same Liouville space (Section~\ref{sec:electron_impact}), shifting
$L \to L + i\Phi(\omega)$:
\begin{equation}
\label{eq:resolvent_broadened}
I(\omega) = \frac{1}{\pi}\operatorname{Re}\,
  \langle\langle \vec{d}^{\,*} \big|\,
  \bigl(\omega\,\hat{I} - L - i\,\Phi(\omega)\bigr)^{-1}
  \,\big| \vec{d}\,\rho_0 \rangle\rangle
\end{equation}

\paragraph{Field-dependent Liouvillian and static profile.}
Under the quasi-static ion approximation the radiator is exposed to a static ionic
microfield $\vec{F}$ at angle $\theta$ to $\vec{B}$.  Both manifold Hamiltonians $H_u$
and $H_l$ then depend on the field magnitude $F$ and orientation $\mu = \cos\theta$
through the Stark perturbation $V_E$ (Section~\ref{sec:full_hamiltonian}), making the
Liouvillian parametrically field-dependent:
\begin{equation}
\label{eq:liouvillian_field}
L(F,\mu) = \frac{1}{\hbar}\Bigl[H_u(F,\mu)\otimes\hat{I}^d
           - \hat{I}\otimes H_l^d(F,\mu)\Bigr],
\qquad H(F,\mu) = H_A + V_E(F,\mu)
\end{equation}
where $H_A$ is the field-free magnetic Hamiltonian (Section~\ref{sec:radiator_H}) and
$V_E(F,\mu)$ is the linear Stark interaction (Section~\ref{sec:full_hamiltonian}).
Equations~(\ref{eq:resolvent_broadened}) and~(\ref{eq:liouvillian_field}) together give
the field-dependent $q$-polarized profile at a given configuration $(F,\mu)$:
\begin{equation}
\label{eq:field_profile}
I_q(\omega,F,\mu) = \frac{1}{\pi}\operatorname{Re}\,
  \langle\langle d_q^*\,\big|\,
  \bigl(\omega\,\hat{I} - L(F,\mu) - i\,\Phi(\omega)\bigr)^{-1}
  \,\big|\,d_q\,\rho_0\rangle\rangle
\end{equation}
The total $q$-polarized profile is then the average over the ionic microfield
distribution $W(F)$ and orientation:
\begin{equation}
\label{eq:microfield_integral}
I_q(\omega) = \frac{1}{2}\int_0^{\infty}\!\int_{-1}^{1}
  W(F)\,I_q(\omega,F,\mu)\;d\mu\;dF
\end{equation}
Because $I_q(\omega,F,\mu) = I_q(\omega,F,-\mu)$ (the Hamiltonian is even in $\mu$
via $F_z = F\mu$ and $F_x = F\sqrt{1-\mu^2}$), the integral over $[-1,1]$
reduces to twice the $[0,1]$ range, discretized by Gauss--Legendre quadrature.
The observed intensity at angle $\alpha$ between the line of sight and $\vec{B}$ is:
\begin{equation}
\label{eq:stokes}
I(\omega,\alpha) = I_\pi(\omega)\sin^2\!\alpha
  + \tfrac{1}{2}\bigl(I_{\sigma^+}(\omega)+I_{\sigma^-}(\omega)\bigr)(1+\cos^2\!\alpha)
\end{equation}
Sections~\ref{sec:radiator_H}--\ref{sec:electron_impact} detail each ingredient
of Eq.~(\ref{eq:field_profile}); Section~\ref{sec:ffm} extends this framework to include
ion dynamics via the FFM.

\subsection{Radiator Hamiltonian in Magnetic Fields}
\label{sec:radiator_H}
For each principal quantum number shell $n$, the atomic Hamiltonian $H_A$ is constructed in the canonical uncoupled hydrogenic basis $|n, l, m_l, s{=}\tfrac{1}{2}, m_s\rangle$ of dimension $2n^2$. Basis states are ordered by an outer loop over $l \in [0, n-1]$, an intermediate loop over $m_l \in [-l, l]$, and an inner loop over $m_s \in \{+\tfrac{1}{2}, -\tfrac{1}{2}\}$. Five physical contributions are assembled into a single matrix \cite{bethesalpeter}:
\begin{equation}
H_A = H_0 + H_{\rm SO} + H_{\rm FS} + H_Z^{(1)} + H_Z^{(2)}
\end{equation}

\subsubsection{Unperturbed Hydrogenic Energy ($H_0$)}
The diagonal unperturbed energy is corrected for the finite nuclear mass:
\begin{equation}
\langle n, l, m_l, m_s | H_0 | n, l, m_l, m_s \rangle = -\frac{Z^2\,\mathrm{Ry}_{\rm red}}{n^2},
\qquad
\mathrm{Ry}_{\rm red} = \mathrm{Ry}_{\infty}\,\frac{M_{\rm nuc}}{m_e + M_{\rm nuc}}
\end{equation}
with $A_{\rm emitter}$ the atomic mass number of the emitting isotope. When \texttt{use\_empirical\_\allowbreak{}data=True} this analytic diagonal is replaced by NIST tabulated energies (Section~\ref{sec:empirical}).

\subsubsection{Spin-Orbit Coupling ($H_{\rm SO}$)}
The spin-orbit interaction $H_{\rm SO} = \xi_{nl}\,\vec{L}\cdot\vec{S}$ couples orbital and spin angular momenta. The coupling constant for $l > 0$ is:
\begin{equation}
\xi_{nl} = \frac{Z^4 \alpha^2\,\mathrm{Ry}_{\infty}}{n^3\,l\,(l + \tfrac{1}{2})\,(l + 1)}
\end{equation}
Using the ladder decomposition $\vec{L}\cdot\vec{S} = L_z S_z + \tfrac{1}{2}(L_+ S_- + L_- S_+)$, the matrix elements within a given $l$-subshell are:
\begin{itemize}
    \item \textbf{Diagonal} ($\Delta m_l = 0,\,\Delta m_s = 0$):
    \begin{equation}
    \langle l, m_l, m_s | H_{\rm SO} | l, m_l, m_s\rangle = \xi_{nl}\,m_l\,m_s
    \end{equation}
    \item \textbf{Off-diagonal} ($\Delta m_l = \pm 1,\,\Delta m_s = \mp 1$): for $s = \tfrac{1}{2}$ the spin ladder factor equals 1 for all allowed transitions:
    \begin{equation}
    \langle l, m_l \pm 1, m_s \mp 1 | H_{\rm SO} | l, m_l, m_s\rangle
    = \tfrac{1}{2}\,\xi_{nl}\sqrt{l(l+1) - m_l(m_l \pm 1)}
    \end{equation}
\end{itemize}

\subsubsection{Fine-Structure Relativistic Corrections ($H_{\rm FS}$)}
To restore the Dirac degeneracy (e.g.\ $2s_{1/2}$ and $2p_{1/2}$), the mass-velocity and Darwin terms are added \cite{bethesalpeter}. These are diagonal in $|n, l, m_l, m_s\rangle$ and independent of $m_l$ and $m_s$:
\begin{equation}
\langle l | H_{\rm FS} | l \rangle =
\begin{cases}
-A_{\rm FS}\!\left(n - \tfrac{3}{4}\right) & l = 0 \\[6pt]
-A_{\rm FS}\!\left(\dfrac{n}{l + \tfrac{1}{2}} - \tfrac{3}{4}\right) & l > 0
\end{cases}
\end{equation}
where $A_{\rm FS} = Z^4\alpha^2\,\mathrm{Ry}_{\infty}/n^4$. The algebraic identity $H_{\rm SO} + H_{\rm FS}$ reproduces the exact Dirac fine-structure eigenvalues for both $j = l \pm \tfrac{1}{2}$.

\subsubsection{Linear Zeeman Effect ($H_Z^{(1)}$)}
\begin{equation}
\langle m_l, m_s | H_Z^{(1)} | m_l, m_s \rangle = \mu_B B\,(m_l + g_s m_s)
\end{equation}
where $\mu_B \approx 5.788\times10^{-5}$\,eV/T and $g_s \approx 2.0023192$.

\subsubsection{Quadratic (Diamagnetic) Zeeman Effect ($H_Z^{(2)}$)}
\begin{equation}
H_Z^{(2)} = \frac{e^2 B^2}{8 m_e}\,r^2\sin^2\theta
\end{equation}
This operator couples states with $\Delta l = 0$ and $\Delta l = \pm 2$. The matrix elements factor into radial and angular parts:
\begin{equation}
\langle l_1, m_l | H_Z^{(2)} | l_2, m_l \rangle
= \left(\frac{e^2 B^2 a_0^2}{8m_e}\right)
  \langle n, l_1 | r^2 | n, l_2 \rangle\;
  \langle l_1, m_l | \sin^2\theta | l_2, m_l \rangle
\end{equation}
\begin{enumerate}
    \item \textbf{Radial integrals}:
    \begin{itemize}
        \item Diagonal ($l_1 = l_2 = l$):
        $\displaystyle\langle n, l | r^2 | n, l \rangle = \frac{n^2}{2Z^2}\bigl[5n^2+1-3l(l+1)\bigr]\;[a_0^2]$
        \item Off-diagonal ($l_2 = l_1 \pm 2$): evaluated numerically via $\int_0^\infty R_{nl_1}\,r^4\,R_{nl_2}\,dr$, cached with \texttt{lru\_cache}.
    \end{itemize}
    \item \textbf{Angular integrals}:
    \begin{itemize}
        \item Diagonal: for $l=0$, $\langle\sin^2\theta\rangle = 2/3$; for $l>0$,
        \begin{equation}
        \langle l, m_l | \sin^2\theta | l, m_l \rangle
        = 1 - \frac{2l^2 + 2l - 1 - 2m_l^2}{(2l-1)(2l+3)}
        \end{equation}
        \item Off-diagonal ($l_2 = l_1 + 2$, $l_{\rm low} = \min(l_1,l_2)$): the constant term in $1 - \cos^2\theta$ vanishes for orthogonal harmonics, leaving:
        \begin{equation}
        \langle l_1, m_l | \sin^2\theta | l_2, m_l \rangle
        = -\sqrt{\frac{\bigl[(l_{\rm low}+1)^2 - m_l^2\bigr]
                       \bigl[(l_{\rm low}+2)^2 - m_l^2\bigr]}
                      {(2l_{\rm low}+1)(2l_{\rm low}+3)^2(2l_{\rm low}+5)}}
        \end{equation}
    \end{itemize}
\end{enumerate}

\subsection{Empirical Field-Free Energies}
\label{sec:empirical}
The analytic diagonal $H_0 = -Z^2\,\mathrm{Ry}_{\mathrm{red}}/n^2$, together with the
spin-orbit and fine-structure terms, reproduces the field-free level positions only
to the accuracy of the hydrogenic Dirac formula. For quantitative line centers,
\texttt{StarkZee} can instead inject \emph{measured} field-free energies through the
\texttt{use\_empirical\_data=True} flag (with the element symbol \texttt{atom}). The
values are read from a tabulated database (\texttt{atomic\_data.load\_levels}) holding
NIST level energies in wavenumbers [cm$^{-1}$], resolved by $(n, l, j)$.
The empirical Hamiltonian is kept in cm$^{-1}$ throughout this code path so that
callers can work entirely in wavenumber units; all Zeeman contributions are
also converted from eV to cm$^{-1}$ before being added, keeping every term on
the same scale.

Because the field-free Hamiltonian is degenerate (the Dirac formula makes
$2s_{1/2}$ and $2p_{1/2}$ coincide), the empirical energies cannot simply be written
on the diagonal of the uncoupled $|n,l,m_l,m_s\rangle$ basis: \texttt{numpy.linalg.eigh}
is free to mix the degenerate $l$ states arbitrarily. \texttt{StarkZee} therefore:
\begin{enumerate}
    \item diagonalizes a degeneracy-broken field-free Hamiltonian (unperturbed energy
    $+$ spin-orbit only, \emph{omitting} the mass-velocity/Darwin term so that $l$
    remains a good label), yielding eigenvectors $V$;
    \item labels each coupled eigenstate $k$ by its dominant orbital component $l$ and
    its total angular momentum $j = l \pm \tfrac{1}{2}$, the branch being set by the
    sign of $\langle \vec{L}\cdot\vec{S}\rangle = (E_k - E_n)/\xi_{nl}$;
    \item assigns the tabulated energy $D_k = E^{\mathrm{emp}}(l, j)$ [cm$^{-1}$] to each eigenstate
    and reconstructs the field-free Hamiltonian as
    \begin{equation}
    H_0^{\mathrm{emp}} = V\,\mathrm{diag}(D)\,V^\dagger \quad [\mathrm{cm}^{-1}].
    \end{equation}
\end{enumerate}
The Zeeman terms ($\mu_B B (m_l + g_s m_s)$ and the diamagnetic correction) are
converted from eV to cm$^{-1}$ and added on top of $H_0^{\mathrm{emp}}$.
Since the tabulated values are absolute level energies (ground state at $0$),
transition wavenumbers follow directly as differences $E_u - E_l$ and reproduce
the observed NIST Lyman/Balmer line centers.

\subsection{Full Electron-Radiator Hamiltonian}
\label{sec:full_hamiltonian}
Under the quasi-static ion approximation, the radiator is subjected to a constant electric microfield $\vec{F}$ at an angle $\theta$ relative to the magnetic field $\vec{B}$. The microfield is decomposed into longitudinal ($F_z = F\cos\theta$) and transverse ($F_x = F\sin\theta$) components. The Stark interaction is
\begin{equation}
V_E = -e\,\vec{r}\cdot\vec{F} = -e(z F_z + x F_x) = -\hat{\boldsymbol{d}}\cdot\vec{F}
\end{equation}
where $\hat{\boldsymbol{d}} = e\vec{r}$ is the electric dipole operator. The full electron-radiator Hamiltonian is
\begin{equation}
\label{eq:full_hamiltonian}
H = H_A + V_E = H_0 + V_\mathrm{SO} + H_Z^{(1)} + H_Z^{(2)} + V_E
\end{equation}
$H_A$ and $V_E$ are assembled into a single matrix and diagonalized \emph{as a whole} by \texttt{numpy.linalg.\allowbreak{}eigh} at every microfield quadrature point (\texttt{solve\_starkzee}). The Stark interaction is \textbf{not} treated perturbatively: there is no expansion in powers of $V_E$. This exact treatment is essential when the Stark energy $eF\langle r\rangle_n$ is comparable to the fine-structure or Zeeman splittings, yielding the Stark-dressed eigenstates $|k\rangle$ and transition frequencies $\omega_k$.

\subsubsection{Stark Perturbation Matrix Elements}
The selection rule for the within-shell Stark coupling is $\Delta n = 0$, $|l_i - l_j| = 1$, $\Delta m_s = 0$. In energy units [eV]:
\begin{equation}
\langle n, l_i, m_{l,i} | V_E | n, l_j, m_{l,j} \rangle
= -\,\delta_{m_{s,i},m_{s,j}}\;
  \langle n, l_i | r | n, l_j \rangle\;
  \bigl\langle l_i, m_{l,i} \big| (z F_z + x F_x)/r \big| l_j, m_{l,j} \bigr\rangle
\end{equation}

\paragraph{Radial matrix element.}
For $l = \max(l_i, l_j)$, Gordon's formula gives the within-shell radial dipole element exactly:
\begin{equation}
\langle n, l | r | n, l-1 \rangle = \frac{3n}{2Z}\sqrt{n^2 - l^2} \quad [a_0]
\end{equation}

\paragraph{Angular matrix elements.}
Writing $z/r = T_0^{(1)}$ and $x/r = (T_{-1}^{(1)} + T_{+1}^{(1)})/\sqrt{2}$:
\begin{itemize}
    \item \textbf{$q=0$} ($\Delta m_l = 0$):
    $\displaystyle
    \langle l, m_l | \cos\theta | l+1, m_l \rangle
    = \sqrt{\dfrac{(l+1)^2 - m_l^2}{(2l+1)(2l+3)}}$
    \item \textbf{$q=+1$} ($\Delta m_l = -1$):
    \begin{equation}
    \langle l_1, m_{l,1} | T_{+1}^{(1)} | l_2, m_{l,2} \rangle =
    \begin{cases}
    +\sqrt{\dfrac{(l - m_{l,2})(l - m_{l,2} - 1)}{2(2l-1)(2l+1)}} & l_1 = l_2 - 1 \\[8pt]
    -\sqrt{\dfrac{(l + m_{l,2})(l + m_{l,2} + 1)}{2(2l-1)(2l+1)}} & l_1 = l_2 + 1
    \end{cases}
    \end{equation}
    \item \textbf{$q=-1$} ($\Delta m_l = +1$):
    \begin{equation}
    \langle l_1, m_{l,1} | T_{-1}^{(1)} | l_2, m_{l,2} \rangle =
    \begin{cases}
    -\sqrt{\dfrac{(l + m_{l,2})(l + m_{l,2} - 1)}{2(2l-1)(2l+1)}} & l_1 = l_2 - 1 \\[8pt]
    +\sqrt{\dfrac{(l - m_{l,2})(l - m_{l,2} + 1)}{2(2l-1)(2l+1)}} & l_1 = l_2 + 1
    \end{cases}
    \end{equation}
\end{itemize}
where $l = \max(l_1, l_2)$ in all cases. The Stark matrices $M_z$ and $M_x$ [eV/(V/m)] are precomputed once per $(n, Z)$ and cached (\texttt{\_stark\_templates}), so that $V_E = F_z M_z + F_x M_x$ is formed at each quadrature point by a simple scalar multiplication.

\subsection{Transition Dipole Operator}
\label{sec:dipole_operator}
The transition dipole matrices $D_q$ ($q \in \{0,\pm1\}$) are built in the uncoupled basis $|n, l, m_l, s, m_s\rangle$. Selection rules require $\Delta l = \pm 1$, $\Delta m_l = q$, and $\Delta m_s = 0$:
\begin{equation}
\langle n_l, l_l, m_{l,l}, m_{s,l} | D_q | n_u, l_u, m_{l,u}, m_{s,u} \rangle
= \delta_{m_{s,l}, m_{s,u}}\;
  \langle n_l, l_l | r | n_u, l_u \rangle\;
  \langle l_l, m_{l,l} | T_q^{(1)} | l_u, m_{l,u} \rangle
\end{equation}

\subsubsection{Angular Dipole Matrix Elements}
The angular factors $\langle l_l, m_{l,l} | T_q^{(1)} | l_u, m_{l,u} \rangle$ follow from the Wigner-Eckart theorem \cite{edmonds}:
\begin{itemize}
    \item For $q=0$ ($\Delta m_l = 0$):
    \begin{equation}
    \langle l, m_l | T_0^{(1)} | l+1, m_l \rangle
    = \sqrt{\frac{(l+1)^2 - m_l^2}{(2l+1)(2l+3)}}
    \end{equation}
    \item For $q=+1$ ($\Delta m_l = -1$):
    \begin{equation}
    \langle l_1, m_{l,1} | T_{+1}^{(1)} | l_2, m_{l,2} \rangle =
    \begin{cases}
    +\sqrt{\dfrac{(l - m_{l,2})(l - m_{l,2} - 1)}{2(2l-1)(2l+1)}} & l_1 = l_2 - 1 \\[8pt]
    -\sqrt{\dfrac{(l + m_{l,2})(l + m_{l,2} + 1)}{2(2l-1)(2l+1)}} & l_1 = l_2 + 1
    \end{cases}
    \end{equation}
    \item For $q=-1$ ($\Delta m_l = +1$):
    \begin{equation}
    \langle l_1, m_{l,1} | T_{-1}^{(1)} | l_2, m_{l,2} \rangle =
    \begin{cases}
    -\sqrt{\dfrac{(l + m_{l,2})(l + m_{l,2} - 1)}{2(2l-1)(2l+1)}} & l_1 = l_2 - 1 \\[8pt]
    +\sqrt{\dfrac{(l - m_{l,2})(l - m_{l,2} + 1)}{2(2l-1)(2l+1)}} & l_1 = l_2 + 1
    \end{cases}
    \end{equation}
\end{itemize}
where $l = \max(l_1, l_2)$.

\subsubsection{Hydrogenic Radial Dipole Integrals (Gordon's Formula)}
For $n_u \neq n_l$, the inter-shell radial elements $\langle n_l, l_l | r | n_u, l_u \rangle$ are evaluated exactly using Gordon's (1929) analytical formula \cite{gordon}:
\begin{multline}
\langle n_l, l_l | r | n_u, l_u \rangle
= \frac{1}{Z}\,\frac{(-1)^{n_l - L}}{4(2L-1)!}
  \sqrt{\frac{(n_u+L)!\,(n_l+L-1)!}{(n_u-L-1)!\,(n_l-L)!}} \\
  \times\frac{(4 n_u n_l)^{L+1}(n_u-n_l)^{n_u+n_l-2L-2}}{(n_u+n_l)^{n_u+n_l}}
  \left[F_1 - \left(\frac{n_u-n_l}{n_u+n_l}\right)^{\!2} F_2\right]
\end{multline}
where $L = \max(l_u, l_l)$, and $F_1$, $F_2$ are terminating Gauss hypergeometric functions:
\begin{align}
F_1 &= {}_2F_1\!\left(-n_u+L+1,\,-n_l+L,\,2L,\,-\frac{4 n_u n_l}{(n_u-n_l)^2}\right) \\
F_2 &= {}_2F_1\!\left(-n_u+L-1,\,-n_l+L,\,2L,\,-\frac{4 n_u n_l}{(n_u-n_l)^2}\right)
\end{align}
Because the first argument in both hypergeometric functions is a non-positive integer, each series terminates as a finite sum:
\begin{equation}
{}_2F_1(-m, b, c, x) = \sum_{k=0}^{m} \frac{(-m)_k\,(b)_k}{(c)_k\,k!}\,x^k
\end{equation}
where $(-m)_k$ is the Pochhammer symbol. This exact evaluation carries no integration cutoff error and is cached for reuse.

\subsection{Diagonalization and Transition Dipole Basis Rotation}
\label{sec:diagonalization}
At each quadrature point $(F, \mu)$ the total Hamiltonians for both manifolds are:
\begin{equation}
H_u(F,\mu) = H_{A,u} + V_{E,u}(F_z, F_x),
\qquad
H_l(F,\mu) = H_{A,l} + V_{E,l}(F_z, F_x)
\end{equation}
Diagonalization via \texttt{numpy.linalg.eigh} gives eigenvectors $V_u$, $V_l$ and eigenvalues $E_u^p$, $E_l^k$:
\begin{equation}
V_u^\dagger H_u V_u = \mathrm{diag}(E_u^p),
\qquad
V_l^\dagger H_l V_l = \mathrm{diag}(E_l^k)
\end{equation}
The transition energy for each dressed-state pair $(p \to k)$ is $E_{pk} = E_u^p - E_l^k$.  The uncoupled dipole matrices $D_q$ are rotated into the mixed eigenstate bases:
\begin{equation}
D'_q = V_l^\dagger\,D_q\,V_u,
\qquad
S_q(p,k) = |D'_q(k,p)|^2
\end{equation}
The weight $S_q(p,k)$ and energy $E_{pk}$ together define one Stark-Dressed Transition (SDT). The static profile at field configuration $(F,\mu)$ — Eq.~(\ref{eq:field_profile}) — is then realized in practice as a sum over all SDTs weighted by the quadrature weight $W_{ij}$:
\begin{equation}
I_q(\omega,F,\mu) = \sum_{p,k} W_{ij}\,S_q(p,k)\,\frac{\gamma_e/\pi}{(\omega - E_{pk})^2 + \gamma_e^2}
\end{equation}
where $\gamma_e$ is the electron-impact half-width of Section~\ref{sec:electron_impact}.

\subsection{Plasma Microfield Distributions}
\label{sec:microfield}
The electric microfield is averaged over a probability distribution $W(F)$. The presence of $\vec{B}$ breaks spherical symmetry, so the integration must be performed over both the field magnitude $F = |\vec{F}|$ and the orientation $\mu = \cos\theta$ between $\vec{F}$ and $\vec{B}$.

\subsubsection{The Hooper Screened Microfield Distribution}
To account for Debye screening of the ion Coulomb fields by surrounding plasma electrons, StarkZee implements the \textbf{Hooper microfield distribution} \cite{hooper}:
\begin{equation}
W(\beta, a) = \frac{2\beta}{\pi} \int_0^\infty y \sin(\beta y)\, T(y, a)\, dy
\end{equation}
where:
\begin{itemize}
    \item $\beta = F / F_0$ is the reduced electric field strength.
    \item $F_0$ is the \textbf{Holtsmark normal field strength}, the Coulomb field of a single perturber at the mean inter-particle distance $r_e$:
    \begin{equation}
    F_0 = \frac{e}{4\pi\varepsilon_0 r_e^2}, \qquad r_e = \left(\frac{3}{4\pi N_e}\right)^{1/3}
    \end{equation}
    \item $a = r_e / \lambda_D$ is the Debye screening parameter, where $\lambda_D$ is the multi-species Debye length:
    \begin{equation}
    \lambda_D = \sqrt{\frac{\varepsilon_0 k_B T_e}{N_e e^2 \bigl(1 + \sum_i X_i Z_i^2\bigr)}}
    \end{equation}
    with $X_i = N_i / N_e$ the fractional ion concentration.
    \item $T(y, a) = \exp\!\bigl(-y^{3/2} S(y,a)\bigr)$ is the screened characteristic function, with Debye screening factor:
    \begin{equation}
    S(y, a) = \left(1 + \frac{1.5\,a^2}{y^2}\right)^{-3/4}
    \end{equation}
\end{itemize}
Setting $a = 0$ yields $S = 1$ and $T = \exp(-y^{3/2})$, recovering the unscreened \textbf{Holtsmark distribution} $W_H(\beta)$.

\subsubsection{The Potekhin Microfield Distribution}
The analytical fits of Potekhin et al.\ (2002) \cite{potekhin} describe the cumulative distribution $Q(\beta)$ and probability density $P(\beta)$ for both screened ($s > 0$) and unscreened ($s = 0$) potentials, for neutral and charged radiators, incorporating the ion-ion coupling parameter $\Gamma$:
\begin{equation}
Q(\beta) = \frac{q_0 \beta^3 - 1.33\,\beta^{9/2} + \beta^6}{q_1 + q_2 \beta^2 + q_3 \beta^3 - \tfrac{1}{3}\beta^{9/2} + \beta^6}
\quad (s=0)
\end{equation}
\begin{equation}
P(\beta) \approx \beta^2 S_N \!\left[ A\,e^{-a\beta^\alpha} + B\,e^{-b\beta^\gamma} + \frac{e^{-\Gamma\beta^{1/2}}}{1 + c\beta^{9/2}} \right]
\quad (s>0)
\end{equation}
where $q_n$, $A, a, \alpha, B, b, \gamma, c$ are tabulated functions of $\Gamma$ and $s$, and $S_N$ is the normalization constant.

\subsubsection{Quadrature Discretization}
The double integral $\iint W(F,\mu)\,dF\,d\mu$ (Eq.~\ref{eq:microfield_integral}) is discretized as follows:
\begin{enumerate}
    \item \textbf{Field magnitude $F$}: A uniform grid of $N_F$ points in reduced field space $\beta_i \in [0, \beta_{\max}]$ (default $\beta_{\max} = 10$):
    \begin{equation}
    F_i = \beta_i F_0, \qquad W_{F,i} = W(\beta_i, a)\,\Delta\beta,
    \qquad \sum_i W_{F,i} = 1
    \end{equation}
    \item \textbf{Field orientation $\mu$}: $N_\mu$-point \textbf{Gauss-Legendre quadrature} over $[0, 1]$. Legendre roots $x_j \in [-1,1]$ and weights $w_j$ are mapped as:
    \begin{equation}
    \mu_j = \frac{x_j + 1}{2}, \qquad W_{\mu,j} = \frac{w_j}{2}
    \end{equation}
    \item \textbf{Combined weight}: $W_{ij} = W_{F,i}\,W_{\mu,j}$. Points with $W_{ij} < 10^{-15}$ are skipped.
\end{enumerate}
At each active point the field decomposes as:
\begin{equation}
F_{z,ij} = F_i\,\mu_j, \qquad F_{x,ij} = F_i\sqrt{1 - \mu_j^2}
\end{equation}
\begin{enumerate}
    \item \textbf{Holtsmark Distribution}: Standard unscreened microfield for weakly coupled plasmas \cite{holtsmark}:
    \begin{equation}
    W_H(\beta) = \frac{2\beta}{\pi} \int_0^\infty y \sin(\beta y) e^{-y^{1.5}} dy
    \end{equation}
    where $\beta = F/F_0$ is the normalized field strength.

    \item \textbf{Hooper Distribution}: Debye-screened microfield distribution \cite{hooper}:
    \begin{equation}
    W_{Hooper}(\beta, a) = \frac{2\beta}{\pi} \int_0^\infty y \sin(\beta y) e^{-y^{1.5} \chi(y, a)} dy
    \end{equation}
    where $\chi(y, a) = (1 + 1.5 a^2 / y^2)^{-0.75}$ is the screening correction fit, and $a = r_e / \lambda_D$ is the ratio of the ion-sphere radius to the Debye length.

    \item \textbf{Potekhin Distribution}: Analytical fits for electric microfield distributions $P(\beta)$ in plasmas \cite{potekhin}, which support both screened and unscreened cases as well as neutral and charged radiator points, incorporating the ion-ion coupling parameter $\Gamma$.
\end{enumerate}

\subsection{Electron Impact Broadening}
\label{sec:electron_impact}
Fast-moving electrons are treated in the \emph{impact} (completed-collision) approximation: each collision is completed in a time short compared to both the inverse linewidth and the mean inter-collision time. In this limit, successive collisions are statistically independent and their net effect on the line shape enters through a damping operator $\Phi(\omega)$ in the Liouville-space resolvent (Eq.~\ref{eq:resolvent_broadened}).

\subsubsection{Liouville-space structure of $\Phi(\omega)$ and connection to $\langle r^2\rangle$}
In the no-quenching approximation (electron collisions do not mix the upper and lower radiating manifolds), the operator $\Phi$ in Liouville space factorizes as a sum of independent contributions from each manifold:
\begin{equation}
\label{eq:phi_liouville}
\Phi(\omega) = \tfrac{1}{2}\bigl(\Phi_u(\omega)\otimes\hat{I}^d + \hat{I}\otimes\Phi_l^d(\omega)\bigr)
\end{equation}
where $\Phi_u$ and $\Phi_l$ are self-energy operators acting within the upper ($n_u$) and lower ($n_l$) manifolds respectively.

The GBK semi-classical model evaluates $\Phi_n$ from the leading long-range term of the electron--radiator interaction, which is the dipole--dipole coupling $V_{ee} \propto r_{\rm atom}^2 / r_e^3$.  Averaging over the Maxwell--Boltzmann distribution of electron impact parameters and velocities, this yields a self-energy proportional to the \emph{operator} $r_n^2$ within shell $n$:
\begin{equation}
\label{eq:phi_r2}
\Phi_n(\omega) = W_0(\omega)\,r_n^2\,\bigl[C_n + G(\omega)\bigr]
\end{equation}
where $r_n^2$ is the mean-square-displacement operator within shell $n$, $W_0$ is the density-temperature prefactor (below), and $G(\omega)$ is the frequency-dependent GBK factor (below).  The $r^2$ dependence is the direct quantum-mechanical trace of the dipole-coupling cross-section: a state with larger electronic extent couples more strongly to a passing electron.

Rotating to the SDT basis via the eigenvectors of $H_n(F,\mu)$ (Section~\ref{sec:diagonalization}) makes $\Phi_n$ approximately diagonal, since the dressed-state splitting $\omega_{kk'}$ is generally large compared to the collision rate.  The diagonal element for upper dressed state $|k\rangle$ is:
\begin{equation}
\label{eq:gamma_k}
\gamma_k^{(u)} = W_0\,\langle k\,|\,r_u^2\,|\,k\rangle\,\bigl[C_{n_u} + G(\Delta E_k)\bigr]
\end{equation}
The resolvent $(\omega\hat{I} - L(F,\mu) - i\Phi(\omega))^{-1}$ then has eigenvalues $(\omega - \omega_k - i\gamma_k)^{-1}$ in the SDT basis, and Eq.~(\ref{eq:field_profile}) reduces directly to the Lorentzian sum of Eq.~(\ref{eq:profile_sum}).

\paragraph{Shell-average approximation.}
Computing the per-SDT matrix element $\langle k | r_u^2 | k \rangle$ exactly requires rotating the $r^2$ operator to each dressed-state basis at every quadrature point (\texttt{electron\_operator=True}).  The default approximation replaces $\langle k | r_u^2 | k \rangle$ by the shell average $\langle r^2 \rangle_{n_u}$, giving a single homogeneous half-width:
\begin{equation}
\gamma_e(\Delta E) = W_0\,\langle r^2 \rangle_{n_u} \bigl[ C_{n_u} + G(\Delta E) \bigr]
\end{equation}
The lower-manifold contribution $\gamma_k^{(l)}$ is retained in principle through Eq.~(\ref{eq:phi_liouville}); in practice, for Balmer lines ($n_l=2$, $\langle r^2\rangle_2 = 12\,a_0^2$) it is a factor ${\sim}3$ smaller than the $n_u=3$ upper term and is dominated by the latter.

\subsubsection{Density-Temperature Prefactor}
\begin{equation}
W_0 = \frac{4\pi}{3}\,N_e\sqrt{\frac{2m_e}{\pi k_B T_e}}
      \left(\frac{\hbar}{m_e}\right)^{\!2}\frac{\hbar}{e}
\quad [\mathrm{eV}\,a_0^{-2}]
\end{equation}
The $\hbar/e$ factor converts rad\,s$^{-1}$ to eV.

\subsubsection{Shell-Averaged Mean-Square Radius}
\begin{equation}
\langle r^2 \rangle_n = \frac{1}{n^2}\sum_{l=0}^{n-1}(2l+1)\,\frac{n^2}{2Z^2}\bigl[5n^2+1-3l(l+1)\bigr]\,a_0^2
\end{equation}
This expression scales as $n^4/Z^2$ and is \emph{not} replaced by an approximate constant.

\subsubsection{Strong-Collision Constant}
Values from Ferri, Peyrusse \& Calisti \cite{ferri}, Table~1:
\begin{equation}
C_n =
\begin{cases}
1.50 & n \le 2 \\
1.00 & n = 3 \\
0.75 & n = 4 \\
0.50 & n = 5 \\
0.40 & n \ge 6
\end{cases}
\end{equation}

\subsubsection{Dynamical GBK Factor}
\begin{equation}
G(\Delta E) = \frac{1}{2}\,E_1\!\left(y(\Delta E)\right),
\qquad
y(\Delta E) = \left(\frac{n^2}{2Z}\right)^{\!2}
              \frac{\Delta E^2 + E_c^2}{\mathrm{Ry}_\infty\,T_e}
\end{equation}
where $\Delta E = E - E_{pk}$ [eV] is the detuning from the transition center, $E_c = \hbar\omega_c$ [eV] is the cutoff energy, $\mathrm{Ry}_\infty \approx 13.606$\,eV is the Rydberg energy ($e^2/2a_0$, the ionization energy of hydrogen), and $E_1(x)$ is the first-order exponential integral function defined for $x > 0$ by:
\begin{equation}
E_1(x) = \int_x^\infty \frac{e^{-t}}{t}\, dt
\end{equation}

\paragraph{Cutoff frequency.}
The collision integral diverges logarithmically at large impact parameters unless truncated at a maximum $\rho_{\max}$, or equivalently at a lower cutoff $\omega_c = v_{\rm th}/\rho_{\max}$ \cite{griem1959}. StarkZee follows the formulation of Ferri, Peyrusse \& Calisti \cite{ferri}:
\begin{equation}
\omega_c = \max\!\left(\omega_p,\;\omega_e,\;\omega_L,\;\omega_{\alpha\alpha'}\right)
\end{equation}
\begin{itemize}
    \item $\omega_p = \sqrt{N_e e^2/(\varepsilon_0 m_e)}$ — \textbf{electron plasma frequency}: Debye screening suppresses interactions at $\rho > \lambda_D = v_{\rm th}/\omega_p$.
    \item $\omega_e = v_{\rm th}/r_e$ — \textbf{configuration-change frequency}: the impact approximation requires each collision to be completed before the surrounding electron configuration changes appreciably. At high density or low $T_e$ this can be shorter than the Debye timescale (see Section~\ref{sec:zest_electron}).
    \item $\omega_L = eB/m_e$ — \textbf{electron Larmor frequency}: in $\vec{B}$, electrons follow helical trajectories, reducing the effective interaction duration.
    \item $\omega_{\alpha\alpha'}$ — \textbf{emitter level-splitting frequency}: zero for hydrogen (degenerate $l$-subshells).
\end{itemize}
The largest frequency dominates, imposing the tightest bound on $\rho_{\max}$.

\paragraph{Frequency-dependent vs.\ resonance-center width.}
When \texttt{frequency\_dependent\_\allowbreak{}width=True} (default), $\gamma_e(\Delta E)$ is evaluated at the actual
detuning of each transition component using the full $E_1(y(\Delta E))$ function.
Setting this flag to \texttt{False} fixes the width at the line-center value $\gamma_e(0)$,
which reduces computation time at the cost of accuracy in the far wings.

\paragraph{Application in the profile.}
Each Stark-dressed transition component $(i \to j, q)$ is broadened by a Lorentzian of half-width $\gamma_e$:
\begin{equation}
L_{ij}(E) = \frac{\gamma_e/\pi}{(E - E_{ij})^2 + \gamma_e^2}
\end{equation}
and the contribution is weighted by the component intensity $|d_q(i,j)|^2$ and the microfield quadrature weight.

\subsubsection{ZEST electron broadening variant (\texttt{electron\_model='zest'})}
\label{sec:zest_electron}

Activating \texttt{electron\_model='zest'} in \texttt{LineProfile.compute\_profile}
switches to the electron broadening model of the ZEST code \cite{calisti2014}.
The formula has the same structure as the default GBK expression but replaces the
full shell-averaged $\langle r^2 \rangle_n$ with the \emph{intra-shell} mean-square
radius:
\begin{equation}
\gamma_e^\text{ZEST}(\Delta E) = W_\mathrm{pref}\,\langle r^2_\mathrm{intra} \rangle_n
\left[ G_n + G_\mathrm{ZEST}(\Delta E) \right]
\end{equation}

\paragraph{Physical justification — $\Delta n = 0$ interaction channels.}
The ZEST paper (Section 2.1 of \cite{calisti2014}) adopts two approximations that
together restrict the electron broadening operator to within-shell matrix elements:

\begin{enumerate}
  \item \textbf{No-quenching approximation}: perturbers may not induce transitions
    between the lower and upper manifolds participating in the radiator emission
    process.
  \item \textbf{$\Delta n = 0$ interaction channels}: ``only states belonging to the
    same Layzer complex [same principal quantum number $n$] may be mixed by Stark
    and/or Zeeman effects.''
\end{enumerate}

A Layzer complex is the set of all $n^2$ states sharing principal quantum number $n$.
These two restrictions reduce the formal sum over all intermediate states in the
broadening operator to within-shell dipole matrix elements only, directly yielding
$\langle r^2_\mathrm{intra}\rangle_n$ in place of the full $\langle r^2\rangle_n$.
The restriction is also automatic from the GBK dynamical factor: an inter-shell
detuning (e.g.\ the $n=3 \to n=2$ gap of $\approx 1.89$ eV at the conditions of
Fig.~2) drives $G(\Delta\omega_\mathrm{inter}) \to 0$ exponentially, so inter-shell
contributions vanish regardless of the explicit approximation.

\paragraph{Intra-shell squared radius.}
Starting from the intra-shell radial element
$\langle n,l|r|n,l{\pm}1\rangle = \tfrac{3n}{2Z}\sqrt{n^2-(l\pm1)^2}$
and angular weight factors $C(l,l+1)=(l+1)/(2l+1)$, $C(l,l-1)=l/(2l+1)$,
the per-$l$ intra-shell sum is
\begin{equation}
r^2_{\mathrm{intra},l} = \frac{9n^2}{4Z^2}\bigl(n^2 - l(l+1) - 1\bigr)
\quad [a_0^2]
\end{equation}
This formula holds for all $l \in [0, n-1]$, including the boundary cases $l=0$
(no $l-1$ neighbor) and $l=n-1$ (no $l+1$ neighbor).
Averaging over the $n^2$ spatial states with $(2l+1)$ weight yields the exact
closed form:
\begin{equation}
\langle r^2_\mathrm{intra}\rangle_n = \frac{9n^2(n^2-1)}{8Z^2}
\quad [a_0^2]
\end{equation}
Numerical values for H ($Z=1$): $n=2$: 13.5\,$a_0^2$; $n=3$: 81\,$a_0^2$;
$n=4$: 270\,$a_0^2$.  These are smaller than the full shell averages
(33, 153, 468\,$a_0^2$) by factors of 0.41, 0.53, and 0.58, respectively.

\paragraph{Maximum wave-number cutoff $\kappa_m$.}
The PPPB model fixes the minimum impact parameter as $\rho_\mathrm{min} = n^2 a_0/(2Z)$ (Bohr orbit)
and uses a combined cutoff $\omega_c = \max(\omega_p, \omega_e, \omega_L)$.  The ZEST convention
instead defines a temperature-dependent wave-number upper limit
\begin{equation}
\label{eq:kappa_m}
\kappa_m = \min\!\left(\frac{Z}{n^2 a_0},\;\frac{Z\sqrt{2m_e k_B T_e}}{\hbar n^2}\right)
\end{equation}
The geometric branch $Z/(n^2 a_0)$ dominates above $T_e \approx \mathrm{Ry}_\infty \approx 13.6$\,eV;
below that the thermal de Broglie branch takes over.  The only cutoff frequency retained in
all ZEST G-functions is the electron plasma frequency $\omega_p$ (no $\omega_e$ or $\omega_L$
terms), so $B$ does not enter the ZEST broadening directly.

\paragraph{G-function choices (\texttt{electron\_model} selector).}
Three alternatives are available for the dynamical factor $G(\Delta\omega)$
in Eq.~(\ref{eq:gamma_k}), all sharing $\langle r^2_\mathrm{intra}\rangle_n$,
the same $G_n = C_n$ strong-collision constants, and $\kappa_m$ as defined above.

\medskip
\textit{\texttt{zest} / \texttt{zest-gbk} (default ZEST).}
Closed-form exponential integral with a $\kappa_m$-derived argument:
\begin{equation}
\label{eq:G_zest_gbk}
G_\mathrm{GBK}(\Delta\omega) = \tfrac{1}{2}\,E_1\!\left(\frac{\Delta\omega^2 + \omega_p^2}{2\,x^2\,\omega_p^2}\right),
\qquad x = \kappa_m \lambda_D
\end{equation}
This has the same $E_1$ structure as the PPPB formula but with $x = \kappa_m\lambda_D$ playing
the role of $\rho_\mathrm{min}/\lambda_D$, and only $\omega_p$ in the additive term.

\medskip
\textit{\texttt{zest-lee}.}
Analytical \emph{min}-blend of the impact ($\Delta\omega \to 0$) and wing ($\Delta\omega \to \infty$) limits:
\begin{align}
G_0 &= \tfrac{1}{2}\!\left[\ln(1+x^2) - \frac{x^2}{1+x^2}\right], \notag\\
G_\infty(\Delta\omega) &= \tfrac{1}{2}\,E_1\!\left(\frac{\Delta\omega^2}{2\,x^2\,\omega_p^2}\right), \notag\\
G_\mathrm{Lee}(\Delta\omega) &= \min\bigl(G_0,\,G_\infty(\Delta\omega)\bigr)
\label{eq:G_lee}
\end{align}
$G_0$ (the line-center plateau) provides the plasma-cutoff floor; the $E_1$ wing term contains
only $\Delta\omega^2$ (no $+\omega_p^2$) because the regularization is already captured by $G_0$.
At large detuning the two forms agree.

\medskip
\textit{\texttt{zest-dufty}.}
Full random-phase-approximation (RPA) numerical integration accounting for Landau damping:
\begin{equation}
\label{eq:G_dufty}
G_\mathrm{Dufty}(\Delta\omega) = \int_{\kappa_\mathrm{min}}^{\kappa_m}
  \frac{e^{-x^2}}{\kappa\,|\varepsilon(\kappa,\Delta\omega)|^2}\,d\kappa,
\qquad x = \frac{|\Delta\omega|}{\kappa\,v_\mathrm{th}},
\quad v_\mathrm{th} = \sqrt{\frac{2k_BT_e}{m_e}}
\end{equation}
where the longitudinal RPA dielectric function is
\begin{equation}
\varepsilon(\kappa,\Delta\omega)
  = 1 + \frac{1}{(\kappa\lambda_D)^2}
    \Bigl[1 - 2x\,\mathcal{D}(x) + i\sqrt{\pi}\,x\,e^{-x^2}\Bigr]
\end{equation}
and $\mathcal{D}(x) = e^{-x^2}\!\int_0^x e^{t^2}\,dt$ is the Dawson function.
The lower limit $\kappa_\mathrm{min} = |\Delta\omega|/(10\,v_\mathrm{th})$ removes the static
($\kappa \to 0$) divergence.  This is the most accurate of the three models but evaluates
one numerical quadrature per frequency point.

\medskip
In practice the three G-functions give nearly identical line cores; differences appear
in the far wings where Landau damping becomes significant.

\subsection{Profile Accumulation}
\label{sec:profile_accumulation}
Each Stark-dressed transition $(p \to k, q)$ is broadened by a Lorentzian of HWHM $\gamma_e(\Delta E)$:
\begin{equation}
L_{pk}(\Delta E) = \frac{1}{\pi}\,\frac{\gamma_e(\Delta E)}{\Delta E^2 + \gamma_e(\Delta E)^2}
\end{equation}
The three polarization profiles are accumulated by vectorized summation over all microfield quadrature points and all SDTs:
\begin{equation}
\label{eq:profile_sum}
P_q(E) = \sum_i W_{F,i} \sum_j W_{\mu,j} \sum_{p,k} S_q(p,k)\,L_{pk}\!\left(E - E_{pk}\right)
\end{equation}
where $q{=}0 \to P_\pi$, $q{=}{-1} \to P_{\sigma^+}$ (blue), $q{=}{+1} \to P_{\sigma^-}$ (red).
This direct summation realizes Eq.~(\ref{eq:microfield_integral}) in practice.

\paragraph{Frequency-dependent vs.\ resonance-center width.}
When \texttt{frequency\_dependent\_\allowbreak{}width=True} (default), $\gamma_e(\Delta E)$ is evaluated at the actual detuning of each $(E, E_{pk})$ pair using the full $E_1(y(\Delta E))$ function. Setting this flag to \texttt{False} fixes the width at the line-center value $\gamma_e(0)$, reducing computation time at the cost of accuracy in the far wings.

\paragraph{Application in the profile.}
Each component is weighted by its SDT intensity $S_q(p,k)$ and the microfield quadrature weight. The observed intensity at angle $\alpha$ to $\vec{B}$ is then:
\begin{equation}
I(\omega,\alpha) = P_\pi(\omega)\sin^2\!\alpha
  + \tfrac{1}{2}\bigl(P_{\sigma^+}(\omega)+P_{\sigma^-}(\omega)\bigr)(1+\cos^2\!\alpha)
\end{equation}

\subsection{Thermal Doppler Broadening}
\label{sec:doppler}

Thermal motion of the radiating ions causes each photon frequency to be Doppler-shifted by $\delta\omega = \omega_0\, v_z/c$. Averaging over a Maxwell--Boltzmann velocity distribution yields a Gaussian line shape with $1/e$ half-width
\begin{equation}
\Delta\omega_D = \omega_0 \sqrt{\frac{2 k_B T_i}{m_\mathrm{ion} c^2}}
\qquad\Longleftrightarrow\qquad
\Delta E_D = E_0 \sqrt{\frac{2 T_i}{m_\mathrm{ion} c^2 / e}}
\end{equation}
where $T_i$ is the ion temperature, $m_\mathrm{ion}$ the emitter mass, and $E_0$ the transition energy in eV. For H Balmer-$\alpha$ at $T_i = 5$\,eV, this gives $\Delta E_D \approx 0.062$\,meV — comparable to the Zeeman splitting at 1\,T and negligible relative to the Stark width at $N_e \sim 10^{23}$\,m$^{-3}$.

\paragraph{Important: Doppler broadening is \emph{not} applied automatically.}
\texttt{calculate\_static\_\allowbreak{}profile()} returns the purely Stark-Zeeman-broadened profile (ion quasi-static + electron impact only). Doppler broadening must be applied \emph{after} as a post-processing step via \texttt{convolutions.py}:
\begin{itemize}
    \item \texttt{apply\_doppler\_broadening(wavelengths\_nm, profile, Ti\_ev, species='H')} — convolves the profile with a Gaussian kernel of the thermal Doppler width.
    \item \texttt{apply\_instrument\_broadening(wavelengths\_nm, profile, fwhm\_nm)} — convolves with a Gaussian instrumental slit function.
\end{itemize}
Both use FFT convolution (\texttt{scipy.fft}) for efficiency. Both operate on a \emph{uniform wavelength grid}; convert energy grids to wavelength before calling. The convolution preserves the total integrated intensity.

The Gaussian kernel applied to the wavelength grid is:
\begin{equation}
K_D(\lambda) = \exp\!\left(-\frac{(\lambda - \lambda_0)^2}{\Delta\lambda_D^2}\right)
\end{equation}
Convolution is performed via the FFT theorem on a wavelength grid edge-padded by $N/2$ on each side to avoid circular-wrap artifacts:
\begin{equation}
P_{q,\mathrm{conv}}(\lambda)
= \mathcal{F}^{-1}\!\left[
    \mathcal{F}(P_{q,\mathrm{pad}})\cdot
    \mathcal{F}\!\bigl(\mathrm{fftshift}(K_{D,\mathrm{pad}})\bigr)
  \right]
\end{equation}
The complete broadening pipeline is therefore:
\begin{equation}
I_\mathrm{obs}(\lambda) = \Bigl[I_\mathrm{SZ}(\lambda) * G_D(\lambda)\Bigr] * G_\mathrm{inst}(\lambda)
\end{equation}
where $I_\mathrm{SZ}$ is the Stark-Zeeman profile from \texttt{static\_profile.py} (via \texttt{line\_profile.py}), $G_D$ is the Doppler Gaussian, and $G_\mathrm{inst}$ the instrumental Gaussian. Both convolutions preserve the total integrated intensity.

\subsection{Frequency Fluctuation Model (FFM)}
\label{sec:ffm}
Dynamic ion motion causes the static microfield to fluctuate over time. In the FFM \cite{talin}, this is modeled as a stationary Markovian process that mixes the Stark-dressed transition components at rate $\nu_i$.

The Stark-Zeeman Hamiltonian $H = H_A + V_E$ is still diagonalized at each field configuration $(F,\mu)$, producing dressed-state frequencies $\omega_k$ and dipole weights $|d_k|^2$ — the Stark-Dressed Transitions (SDTs). This step is purely static: $\omega_k$ and $|d_k|^2$ depend only on the instantaneous ion field. Fast electron collisions add a homogeneous Lorentzian half-width $\gamma_k$ to each SDT (the GBK width from Section~\ref{sec:electron_impact}), acting on a timescale short enough that the ion configuration does not change during a single collision. Ion dynamics enter exclusively through $\nu_i$, estimated as the inverse time for an ion to cross the mean ion spacing at its thermal speed:
\begin{equation}
\nu_i = \frac{v_{th}}{r_i}, \qquad v_{th} = \sqrt{\frac{2 k_B T_i}{m_i}}, \qquad r_i = \left(\frac{3}{4\pi N_i}\right)^{1/3}
\end{equation}
The dynamic line profile is then given by the Sherman-Morrison form:
\begin{equation}
I(\omega) = \frac{r^2}{\pi} \text{Re} \frac{S(\omega)}{1 - \nu_i S(\omega)}
\end{equation}
with the static propagator sum:
\begin{equation}
S(\omega) = \sum_k \frac{p_k}{\nu_i + \gamma_k + i(\omega - \omega_k)}
\end{equation}
where $p_k = |d_k|^2 / r^2$ are the normalized SDT weights and $r^2 = \sum_k |d_k|^2$.
The limit $\nu_i \to 0$ recovers the static profile; $\nu_i \to \infty$ gives a single
Lorentzian (motional narrowing).

\section{Physical Features and Validation}

\subsection{Quadratic Zeeman polarization wings}
At $B \geq 500$\,T the diagonal QZ shifts within $n=3$ differ by up to 5\,meV,
separating the 3p($m_l = \pm 1$)$\to$2s transitions (which carry $\approx 21\%$ of
the H$\alpha$ oscillator strength and exhibit the \emph{largest} n=3 QZ shift,
$+8.87$\,meV at $B=1000$\,T) from the dominant 3d$\to$2p cluster by
$\approx 7$\,meV. The resulting ``wings'' appear only in the
\texttt{quadratic\_zeeman=True} profiles and are validated in
\texttt{test\_12\_halpha\_qz\_wings.py}.

\subsection{\texorpdfstring{$\pm 2\mu_B B$}{±2 mu_B B} Stark-Zeeman satellite features}
A transverse microfield $F_x$ couples adjacent-$m_l$ states within the same
principal shell via $\Delta l = \pm 1,\, \Delta m_l = \pm 1$ matrix elements.
The upper eigenstate near $E_{0,n} + 2\mu_B B$ (dominated by $|nd,\, m_l{=}2\rangle$)
acquires a small $|np,\, m_l{=}1\rangle$ admixture $\beta$, producing a
$\sigma^+$ transition to the zero-Zeeman lower state at photon energy
$E_0 + 2\mu_B B$.

\paragraph{H$\beta$ (n=4$\to$2).}
The satellite is a \emph{distinct peak} at $+117$\,meV ($\approx 2\%$ of the
main $\sigma^+$ intensity), because $n=4$ has both $|4d,\, m_l{=}2\rangle$
and $|4f,\, m_l{=}2\rangle$ \emph{degenerate} at $+2\mu_B B$, greatly
amplifying Stark mixing. The satellite is visible in the Balmer-series profile
at $\approx 465$\,nm and $\approx 509$\,nm for $B=1000$\,T.

\paragraph{H$\alpha$ (n=3$\to$2).}
Only $|3d,\, m_l{=}2\rangle$ sits at $+2\mu_B B$ (no $l=3$ sub-states exist
for $n=3$). The satellite amplitude ($\approx 0.07\%$ of main peak) is
comparable to the Lorentzian tail of the $\sigma^+$ main peak at that
position, so no distinct local maximum appears at 618\,nm / 700\,nm.
This result is converged with respect to the microfield grid size
(\texttt{num\_f} = 30, 60, 100 give identical profiles) and is validated in
\texttt{test\_13\_stark\_zeeman\_satellites.py}.

\paragraph{$B$-dependence and observability.}
The mixing coefficient $\beta = F_x\,\langle np|r|nd\rangle/({\mu_B B})$
scales as $B^{-1}$, so $\beta$ is \emph{larger} at lower $B$. However, the
satellite position $2\mu_B B$ also shrinks with $B$. At $B = 100$\,T,
$2\mu_B B = 11.6$\,meV falls inside the Stark-broadened $\sigma^+$ cluster
($\sim10$\,meV width for $n=4$, $N_e = 10^{23}$\,m$^{-3}$). The satellite
merges into the cluster tail: the peak/valley contrast in the satellite
sub-window is $\approx 1.000$ (indistinguishable), versus $\approx 2.26$ at
$B = 1000$\,T where the satellite is clearly separated. \emph{High $B$ is
therefore required for the satellite to appear as a resolved, separated peak.}

\section{Physical and Numerical Simplifications}
To achieve absolute self-sufficiency, maintainability, and high execution speeds, the following simplifications were incorporated:
\begin{enumerate}
    \item \textbf{Analytical Hydrogenic Basis}:
    Instead of calling external commercial databases or complex atomic structure codes (e.g., Cowan or FAC), SZLS implements exact analytical hydrogenic radial wavefunctions \cite{bethesalpeter}:
    \begin{equation}
    R_{nl}(r) = \sqrt{\left(\frac{2Z}{n}\right)^3 \frac{(n-l-1)!}{2n(n+l)!}} e^{-Zr/n} \left(\frac{2Zr}{n}\right)^l L_{n-l-1}^{2l+1}\left(\frac{2Zr}{n}\right)
    \end{equation}
    This provides self-contained matrix elements for $r$ and $r^2$ calculations.
    
    \item \textbf{Resonance-center impact approximation (optional)}:
    By default (\texttt{frequency\_de\-pendent\_width=True}) the GBK electron width $\gamma_e(\Delta E)$ is evaluated at the actual detuning of each transition component. Setting \texttt{frequency\_de\-pendent\_width=False} fixes the width at its line-center value $\gamma_e(0)$, avoiding repeated evaluations of $E_1(y)$ at every wavelength point. This is a useful approximation when the profile wings are not of primary interest.
    
    \item \textbf{Gauss-Legendre Angle Quadrature}:
    Integrating over the field orientation is performed using a $N_\mu$-point Gauss-Legendre quadrature over $\mu = \cos\theta \in [0, 1]$; see Section~\ref{sec:microfield} for the full discretization scheme.
\end{enumerate}


\section{Built-in Reference Models}
\label{sec:models}

The \texttt{starkzee.models} package provides five independent lineshape models that share a common call signature and serve as benchmarks against StarkZee's fully coupled solver.  They are grouped into tabulated models (reading precomputed databases) and analytical models.

\subsection{Tabulated Models}

\subsubsection{Stehlé (MMM) — \texttt{stehle}}

The Stehlé model reads precomputed Stark lineshapes from the Model Microfield Method (MMM) database \cite{stehle} for hydrogen Balmer and Paschen transitions over a grid of $N_e$ and $T_e$ values.  The tables contain the \emph{unmagnetized} ($B = 0$) Stark + fine-structure profile as a function of a reduced detuning $\Delta\omega/F_0$ (Holtsmark normal field units).  The pipeline is:
\begin{enumerate}
    \item Bilinear interpolation over the $(N_e, T_e)$ table grid (with a Fortran-style three-point hyperbolic interpolator for the detuning axis).
    \item Thermal Doppler broadening via FFT convolution.
    \item Normal Zeeman splitting applied \emph{after} the table: the whole Stark+Doppler profile is rigidly copied to $\omega \pm \omega_L$ (Larmor frequency) with angle-dependent $\pi/\sigma$ weights:
        \begin{equation}
        I(\theta) = \sin^2\!\theta\,I_\pi + \tfrac{1+\cos^2\!\theta}{2}\,(I_{\sigma+} + I_{\sigma-})
        \end{equation}
\end{enumerate}
The Stark and Zeeman effects are therefore treated as \textbf{separable}.  This is valid only when the Zeeman splitting $\mu_B B$ is much smaller than the Stark width.  Stehlé covers arbitrary $(n_u, n_l)$ and species with no restriction on $B$-field value (it is simply ignored in the table lookup).

\subsubsection{Rosato — \texttt{rosato}}

The Rosato database \cite{rosato} solves the Stark-Zeeman problem \emph{jointly} and tabulates the result with $B$ as an explicit table axis ($B \in \{0, 1, 2, 2.5, 3, 5\}$ T).  The pipeline is:
\begin{enumerate}
    \item Interpolation over the $(N_e, T_e, B)$ grid.  The $B$-interpolation is \emph{scaled}: before blending two bracketing profiles at $B_0$ and $B_1$, each profile's detuning axis is stretched by $B_\mathrm{node}/B$.  This exploits the fact that Zeeman splitting scales linearly with $B$, aligning the $\sigma$ components before interpolating.
    \item Angle dependence from \emph{two real tables} (parallel and perpendicular to $\vec{B}$), blended as $I = I_\parallel\cos^2\theta + I_\perp\sin^2\theta$.  Both tables carry the correct $\pi/\sigma$ lineshape and width — the angle enters as a physical interpolation, not just an amplitude weight.
    \item Thermal Doppler broadening via FFT convolution.
\end{enumerate}
Rosato is restricted to deuterium Balmer lines and the table coverage $N_e \in [10^{13}, 10^{16}]$ cm$^{-3}$, $T_e \in [0.316, 31.6]$ eV, $B \le 5$ T.



\subsection{Analytical Models}

\subsubsection{Lomanowski — \texttt{lomanowski}}

Polynomial fits for the Stark FWHM from Lomanowski \textit{et al.} \cite{lomanowski}:
\begin{equation}
\Delta\lambda_S = c_{n_u n_l}\,N_e^{a_{n_u n_l}}\,T_e^{-b_{n_u n_l}} \quad [\mathrm{nm}]
\end{equation}
Coefficients $(a, b, c)$ are tabulated for H/D Balmer and Paschen lines up to $n_u = 9$.  A Voigt profile combines this Stark Lorentzian with a Doppler Gaussian.  No magnetic field treatment.

\subsubsection{Parameterized Stehlé — \texttt{stehle\_param}}

An analytical parametric fit to the Stehlé tables, providing a fast closed-form approximation to the field-free Stark profile without reading the NetCDF database.

\subsubsection{Voigt — \texttt{voigt}}

A pseudo-Voigt profile using the Griem $\alpha_{12}$ Stark half-width \cite{griem} combined with thermal Doppler broadening.  Serves as a quick estimate when only order-of-magnitude Stark width accuracy is required.

\subsection{Comparison with StarkZee}

\begin{table}[H]
\centering\small
\begin{tabularx}{\textwidth}{@{}lXXX@{}}
\toprule
\textbf{Feature} & \textbf{StarkZee} & \textbf{Stehlé} & \textbf{Rosato} \\
\midrule
Magnetic field      & Full simultaneous diagonalization of $H = H_A + V_E$
                    & $B = 0$ in tables; added as rigid triplet after
                    & $B$ is a table axis \\[2pt]
Fine structure      & Yes — spin-orbit, MV, Darwin
                    & No (degenerate hydrogenic)
                    & In tables \\[2pt]
Electron broadening & GBK semi-classical, frequency-dependent
                    & Unified theory (more rigorous far wings)
                    & In tables \\[2pt]
Ion dynamics        & Quasi-static + optional FFM
                    & Static + some ion-dynamics treatment
                    & In tables \\[2pt]
Microfield          & Hooper screened & Holtsmark/Hooper variant & Holtsmark variant \\[2pt]
$B$-treatment       & Intrinsic to Hamiltonian & External convolution & $B$-scaled interpolation \\
\bottomrule
\end{tabularx}
\end{table}

The most important physical distinction at $B \ne 0$: StarkZee diagonalizes $H = H_A + V_E$ \emph{simultaneously} for all field strengths.  When $\mu_B B \sim 3n\,e\,a_0\,F$ (Zeeman and Stark splitting comparable) the eigenstates are genuine Stark-Zeeman hybrids — neither pure Zeeman nor pure Stark states.  No post-processing convolution can reproduce this mixing.  The separable approximation (Stehlé) and scaled-interpolation approach (Rosato) both become inaccurate in this regime.

\section{Code Structure and Modular Architecture}
The \textbf{StarkZee} package (\texttt{starkzee/}) is organized as a modular Python library:
\begin{itemize}
    \item \textbf{\texttt{starkzee/line\_profile.py}}: High-level \texttt{LineProfile} class. Manages spectral grids, unit conversions, Doppler broadening, and exposes all results as attributes. Entry point for most users.
    \item \textbf{\texttt{starkzee/static\_profile.py}}: Core solver. \texttt{calculate\_static\_profile()} runs the microfield quadrature loop and accumulates the three polarization profiles $P_\pi$, $P_{\sigma^+}$, $P_{\sigma^-}$.
    \item \textbf{\texttt{starkzee/radiator.py}}: Constructs $H_A$ — the atomic Hamiltonian including unperturbed energy, spin-orbit, fine-structure, and linear and quadratic Zeeman terms — and builds the dipole operator matrices using analytical angular matrix elements.  $H_A$ is \emph{not} diagonalized here; the full $H = H_A + V_E$ is diagonalized in \texttt{solve\_starkzee}.
    \item \textbf{\texttt{starkzee/ffm.py}}: Implements dynamic ion-field fluctuations via the Frequency Fluctuation Model, supporting both an analytical Sherman-Morrison solver and a full matrix inversion.
    \item \textbf{\texttt{starkzee/microfield.py}}: Generates screened Hooper and unscreened Holtsmark microfield distributions; supports multi-species Debye lengths and custom tabulated distributions.
    \item \textbf{\texttt{starkzee/broadening.py}}: GBK electron-impact half-width $\gamma_e(\Delta E)$; computes plasma, Larmor, and configuration-change cutoff frequencies.
    \item \textbf{\texttt{starkzee/convolutions.py}}: FFT-based post-processing convolutions operating on uniform wavelength grids:
    \begin{equation}
    I_\mathrm{conv}(\lambda) = \mathcal{F}^{-1}\!\left\{\mathcal{F}[I(\lambda)]\cdot\mathcal{F}[K_\mathrm{Doppler}(\lambda)]\cdot\mathcal{F}[K_\mathrm{slit}(\lambda)]\right\}
    \end{equation}
    \item \textbf{\texttt{starkzee/models/}}: Comparison lineshape models (voigt, lomanowski, stehle\_param, stehle, rosato) sharing a common signature, centered on NIST air wavelengths for each transition \cite{nist}.
    \item \textbf{\texttt{starkzee/atomic\_loader.py}}: Placeholder for future multi-electron species support.
\end{itemize}

\begin{thebibliography}{9}
\bibitem{ferri}
S. Ferri, O. Peyrusse, A. Calisti, et al., \textit{Stark-Zeeman line-shape model for multi-electron ions in hot dense plasmas}, Matter and Radiation at Extremes, 7, 015901 (2022). DOI: 10.1063/5.0058552
\bibitem{baranger}
M. Baranger, \textit{Simplified Quantum-Mechanical Theory of Pressure Broadening}, Phys. Rev. 111, 481 (1958); \textit{Problem of Overlapping Lines in the Theory of Pressure Broadening}, Phys. Rev. 112, 855 (1958).
\bibitem{bethesalpeter}
H.A. Bethe, E.E. Salpeter, \textit{Quantum Mechanics of One- and Two-Electron Atoms}, Springer-Verlag, Berlin (1957).
\bibitem{gordon}
W. Gordon, \textit{Zur Berechnung der Matrizen beim Wasserstoffatom}, Ann. Phys. 394, 1031 (1929).
\bibitem{edmonds}
A.R. Edmonds, \textit{Angular Momentum in Quantum Mechanics}, Princeton University Press (1957).
\bibitem{talin}
B. Talin, A. Calisti, L. Godbert, R. Stamm, R.W. Lee, \textit{Frequency fluctuation Model for line-shape calculations in Plasma Spectroscopy}, Phys. Rev. A 51, 1918 (1995).
\bibitem{griem}
H.R. Griem, \textit{Principles of Plasma Spectroscopy}, Cambridge University Press (1997).
\bibitem{griem1959}
H.R. Griem, A.C. Kolb, K.Y. Shen, \textit{Stark Broadening of Hydrogen Lines in a Plasma}, Phys.\ Rev.\ 116, 4 (1959).
\bibitem{griembaranger}
H.R. Griem, M. Baranger, A.C. Kolb, G. Oertel, \textit{Stark Broadening of Neutral Helium Lines in a Plasma}, Phys. Rev. 125, 177 (1962).
\bibitem{holtsmark}
J. Holtsmark, \textit{\"Uber die Verbreiterung von Spektrallinien}, Ann. Phys. 58, 577 (1919).
\bibitem{hooper}
C.F. Hooper, Jr., \textit{Low-Frequency Component of the Electric Microfield in a Plasma}, Phys. Rev. 165, 215 (1968).
\bibitem{potekhin}
A.Y. Potekhin, G. Chabrier, D. Gilles, \textit{Analytic fits for electric microfield distributions in plasmas}, Phys. Rev. E 65, 036412 (2002).
\bibitem{nist}
A. Kramida, Yu. Ralchenko, J. Reader, NIST ASD Team, \textit{NIST Atomic Spectra Database (ver. 5.11)}, National Institute of Standards and Technology, Gaithersburg, MD (2023). DOI: 10.18434/T4W30F
\bibitem{stehle}
C. Stehlé, R. Hutcheon, \textit{Extensive tabulations of Stark broadened hydrogen line profiles}, Astron. Astrophys. Suppl. Ser. 140, 93 (1999).
\bibitem{rosato}
J. Rosato, H. Capes, R. Stamm, \textit{Influence of ion dynamics and multiple body effects on hydrogen lines in magnetized fusion plasmas}, Phys. Rev. E 79, 046408 (2009).
\bibitem{lomanowski}
B. A. Lomanowski et al., \textit{Inferring divertor plasma properties from hydrogen Balmer and Paschen series spectroscopy in JET-ILW}, Nucl. Fusion 55, 123028 (2015).
\bibitem{calisti2014}
A. Calisti, S. Ferri, B. Talin, \textit{ZEST: A fast and accurate spectral line-shape code}, Journal of Physics B: Atomic, Molecular and Optical Physics 47, 175701 (2014).
\end{thebibliography}

% ============================================================
\newpage
\appendix
% ============================================================

\section{Algorithmic Overview of the Calculation Pipeline}
\label{app:pipeline}

When a user invokes \texttt{LineProfile.compute\_profile(grid, grid\_type)}, the \texttt{StarkZee} package performs a series of nested calculations to solve the coupled Stark-Zeeman line-shape problem. The numerical delegation chain is structured as follows:

\begin{enumerate}
    \item \textbf{\texttt{LineProfile.compute\_profile()}}: Manages the spectral coordinate axes, performs grid unit conversions, and forwards the arguments to the profile calculations.
    \item \textbf{\texttt{static\_profile.calculate\_static\_profile()}}: Generates the plasma microfield grid and weights, computes the Gauss-Legendre quadrature points for the microfield orientation angle $\theta$, and runs the main double-integration loop.
    \item \textbf{\texttt{solve\_starkzee()}}: Assembles the full electron-radiator Hamiltonian $H = H_A + V_E$ (Eq.~\ref{eq:full_hamiltonian}) and diagonalizes it as a single matrix with \texttt{numpy.linalg.eigh}.  $V_E$ is \emph{not} applied perturbatively.
    \item \textbf{\texttt{radiator.build\_hamiltonian()}}: Builds $H_A$ — the field-free plus magnetic Hamiltonian (unperturbed energy, spin-orbit, fine structure, linear and quadratic Zeeman) in the uncoupled $|n,l,m_l,m_s\rangle$ basis.
    \item \textbf{\texttt{static\_profile.\_stark\_templates()}}: Returns the field-independent Stark coupling matrices $M_z$ and $M_x$ [eV/(V/m)], cached per $(n,Z)$, so that $V_E = F_z M_z + F_x M_x$ is formed at each quadrature point without rebuilding the matrix loop.
    \item \textbf{Transition Dipole Rotation}: Rotates the transition dipole operator matrices $D_q$ to the mixed eigenstate bases.
    \item \textbf{Electron-Impact Broadening}: Evaluates the frequency-dependent electron Stark half-width $\gamma_e(\Delta E)$ via the Griem-Baranger-Kolb (GBK) model and accumulates Lorentzian profiles.
    \item \textbf{\texttt{convolutions.\allowbreak{}apply\_doppler\_\allowbreak{}broadening()}}: Performs a post-processing Fast Fourier Transform (FFT) convolution to apply thermal Doppler broadening.
\end{enumerate}

% ============================================================

\section{Physical Approximations and Domains of Validity}
\label{app:approx}

\begin{table}[H]
\centering\small
\begin{tabular}{@{}p{3.1cm}p{3.4cm}p{3.4cm}p{3.6cm}@{}}
\toprule
\textbf{Approximation} & \textbf{Mathematical form} & \textbf{Domain of validity} & \textbf{Failure mode} \\
\midrule
\textbf{Within-Shell Isolation} ($\Delta n = 0$)
  & $V_E = \mathrm{diag}_n(V_E)$, no $n \to n\!\pm\!1$ coupling
  & Stark shift $\ll$ shell spacing $2Z^2\mathrm{Ry}/n^3$
  & Quadratic Stark; Inglis-Teller merging at $N_e \gtrsim 10^{24}$\,m$^{-3}$. \\[4pt]
\textbf{Quasi-Static Ion Microfields}
  & $\vec{F}_\mathrm{ion} = \mathrm{const}$
  & Fluctuation rate $\nu_i \ll$ Stark width $\Delta E_S$
  & Ion dynamics / motional narrowing; use FFM. \\[4pt]
\textbf{Semi-Classical Electron Impact}
  & Lorentzian with GBK width $\gamma_e$
  & $\lambda_{dB} \ll$ impact parameter $\rho$
  & Quantum close-coupling needed for cold dense plasmas. \\[4pt]
\textbf{Dipole Approximation}
  & $V_\mathrm{rad} \propto \vec{r}\cdot\vec{E}$
  & $\lambda \gg \langle r\rangle_n \sim n^2 a_0/Z$
  & Quadrupole/octupole for high-$n$ Rydberg states. \\
\bottomrule
\end{tabular}
\caption{Key physical approximations, their domains of validity, and failure modes.}
\end{table}

% ============================================================

\section{Physics Explanation of Use Cases and Operational Limits}
\label{app:usecases}

\subsection{The High-Magnetic-Field Regime}
In laboratory magnetic confinement fusion devices (tokamaks, stellarators), the magnetic field reaches $B \approx 1$--$10$\,T:
\begin{itemize}
    \item The Zeeman splitting $\Delta E_Z \approx 0.1$\,meV is comparable to the fine-structure splitting of lower shells, requiring the full coupled Hamiltonian diagonalization.
    \item Typical microfield strengths $F \sim 10^5$--$10^6$\,V/m represent a weak-Stark regime where Stark broadening acts as a symmetric perturbation on the Zeeman triplet structure.
\end{itemize}

For astrophysical compact objects (magnetized white dwarfs with $B \sim 10^3$--$10^5$\,T, neutron stars with $B \sim 10^8$\,T):
\begin{itemize}
    \item The quadratic Zeeman term $H_Z^{(2)} \propto B^2$ dominates, causing significant blue-shifting and highly asymmetric splitting. The \texttt{quadratic\_zeeman=True} flag enables exact numerical treatment of the $\Delta l = \pm 2$ coupling (see Section~\ref{sec:radiator_H}).
    \item \texttt{StarkZee}'s exact numerical computation of $\langle n, l_1 | r^2 | n, l_2\rangle$ avoids the geometric-mean overestimation of up to 41\% for $n=5$.
\end{itemize}

\subsection{Density Limits}
\begin{itemize}
    \item \textbf{Low density} ($N_e \lesssim 10^{18}$\,m$^{-3}$): The profile is dominated by thermal Doppler broadening; microfield weights converge to a delta function at $F=0$.
    \item \textbf{High density} ($N_e \gtrsim 10^{24}$\,m$^{-3}$): The inter-particle spacing $r_e$ approaches the atomic radius $\langle r\rangle_n$. The static and $\Delta n = 0$ approximations break down as the Stark shift exceeds the Rydberg shell spacing (Inglis-Teller limit).
\end{itemize}

\end{document}
